\documentclass[10pt,a4paper]{amsart}
\usepackage[margin=19mm]{geometry}
\usepackage{amsmath,amssymb,mathtools}
\usepackage[scr=rsfs,cal=euler]{mathalfa}
\usepackage{xfrac}
\usepackage[unicode,hidelinks,bookmarks=false]{hyperref}
\newcommand{\ie}{i.e.\ }
\newcommand{\cf}{cf.\ }
\newcommand{\resp}{respectively\ }
\newcommand{\Subsection}{\subsection}
\providecommand{\nobreakdash}{\nobreak\hbox{-}}
\setlength{\emergencystretch}{2em}
\multlinegap0pt
\usepackage{bm}
\usepackage{bbm}
\usepackage{dsfont}
\usepackage{xspace}
\usepackage{cancel}
\usepackage{mathtools}
\usepackage{amsthm}
\makeatletter
\@ifundefined{theorem}{\newtheorem{theorem}{Theorem}}{}
\@ifundefined{corollary}{\newtheorem{corollary}[theorem]{Corollary}}{}
\@ifundefined{proposition}{\newtheorem{proposition}[theorem]{Proposition}}{}
\makeatother
\newcommand{\F}{\mathbb{F}}
\newcommand{\legl}[2]{\left( \frac{#1}{#2} \right)_\ell}
\title[Nonvanishing of $L$--functions associated to fixed order cha]{Nonvanishing of $L$--functions associated to fixed order characters over function fields}
\author{Chantal David, Alexandra Florea, Matilde Lalin}
\date{Source: arXiv:2506.07815v1 (CC BY 4.0)}
\begin{document}
\maketitle

\noindent\emph{Source and licence.} Nonvanishing of $L$--functions associated to fixed order characters over function fields. Version arXiv:2506.07815v1 (CC BY 4.0), distributed under the Creative Commons Attribution 4.0 International licence, as recorded in the arXiv licence metadata for this exact version and shown on the abstract page https://arxiv.org/abs/2506.07815v1. Full rights notice: licenses/arxiv-2506.07815-CC-BY-4.0.txt.\par\smallskip

\noindent\textit{Editorial scope.} Counted unit: the Proposition whose source label is \texttt{None} in the article (source0.tex lines 1155--1164, source label \texttt{None}), delivered together with the article's own complete original proof of it (source0.tex lines 1167--1251, one \texttt{proof} environment of 85 lines). No mathematics is written, repaired or re-derived: statement and proof are the article's own text line for line. Required context retained uncounted: thm:gettingridofa (lines 874--886); lov (lines 1005--1012); sumv (lines 1045--1051); cor\_lov (lines 1143--1149). Every definition, display and label of those lines is retained unchanged. Omitted from this package and not redistributed: 1--873, 887--1004, 1013--1044, 1052--1142, 1150--1154, 1165--1166, 1252--2136. Nothing inside a retained excerpt is omitted or altered: every retained body is delivered exactly as the article's lines are written, so no omission receipt of any kind accompanies this package. Citations: the retained excerpts carry no citation command at all, so no bibliography entry of the article is transcribed and no reference is redistributed. Reference adaptation: none was needed and none was made; every reference inside the delivered unit resolves inside it, so no unresolved reference or citation is produced. Printed slips kept literal: no printed word, symbol, hypothesis or number of the counted statement or of its proof was altered. Layout and class: the article's own class is \texttt{amsart} with packages including fullpage, bm, xy, tikz, amsmath, amsthm, amscd, amssymb, amsfonts, latexsym, mathrsfs, bbm, dsfont, relsize; the wrapper is the lane's pinned \texttt{amsart} editorial wrapper at 10pt on A4, which restates the theorem-like environments the retained text needs and adds the article's own macro definitions that the retained text needs (\texttt{\textbackslash F}, \texttt{\textbackslash legl}), with extra wrapper packages (bm, bbm, dsfont, xspace, cancel, mathtools). The delivered numbering is the unit's own. Assets: no retained span contains a figure environment, a graphics-inclusion command, a PDF-page inclusion, a table or a TikZ picture; no image file is copied into this package, the package declares no asset dependency and no image file is redistributed with it. Licence: version arXiv:2506.07815v1 of this article is distributed under the Creative Commons Attribution 4.0 International licence, as recorded in the arXiv licence metadata for this exact version and shown on the abstract page https://arxiv.org/abs/2506.07815v1. A full rights notice is delivered at licenses/arxiv-2506.07815-CC-BY-4.0.txt. Characters: every retained excerpt is pure ASCII, so the delivered engine, XeLaTeX with its default Latin Modern text font, reports no missing character.
\begin{theorem} \label{thm:gettingridofa} Let $a, r \in \F_q[t]$, where $a$ is monic  square-free and $(r,a)=1$. We write $r = r_1 r_2^2 \cdots r_{\ell-1}^{\ell-1}r_\ell^\ell$, where the $r_j$ are square-free and pairwise coprime
for $j=1,\dots,\ell-1$, and let $r_\ell^*$ be the product of the primes dividing $r_\ell$, but not $r_1\cdots r_{\ell-1}$.
Then 
\begin{align*}
  \psi_r^{(i)}(r, a,u) =&
 G_\ell(r,a) u^{\deg(a)}\sum_{\substack{E\in \mathcal{M}\\ E \mid r_\ell^* }}\mu(E)G_\ell(a^{\ell-2}r_1\cdots r_{\ell-1}^{\ell-1},E)u^{\deg(E)}\\&\times
  \prod_{\pi \mid aEr_1\cdots r_{\ell-1}}
\left(1 - q^{(\ell-1)\deg(\pi)} \legl{-1}{\pi}^{\nu_\pi(aEr_1\cdots r_{\ell-2}r_{\ell-1})+1}u^{\ell \deg(\pi)}\right)^{-1}\nonumber \\
&\times  \sum_{\substack{d_{\ell-2} \mid aEr_{\ell-2}\\ d_1\mid r_1, \dots, d_{\ell-3}\mid r_{\ell-3}}} \mu(d_1d_2\cdots d_{\ell-2} ) q^{\deg(d_1d_2^2\cdots d_{\ell-2}^{\ell-2})}  u^{\deg(d_1^2d_2^3\cdots d_{\ell-2}^{\ell-1}  )} \prod_{1\leq k< j \leq \ell-2} \overline{\legl{d_k}{d_j}^{k(j+1)}}\\
& \times \prod_{j=1}^{\ell-2} G_{\ell,j+1} \left( \frac{a^{\ell-2}E^{\ell-2}r_1 r_2^2 \cdots r_{\ell-1}^{\ell-1}}{d_{1}d_2\cdots d_{\ell-2} ^{\ell-2}} , d_{j} \right)\\
&  \times \psi^{(i-\deg(a)-\deg(E)- \deg(d_1^2d_2^3\cdots d_{\ell-2}^{\ell-1} ))} \left( \frac{a^{\ell-2}E^{\ell-2}r_1 r_2^2 \cdots r_{\ell-1}^{\ell-1}}{\prod_{j=1}^{\ell-2}d_j^{2j+2-\ell}}, u \right).
\end{align*}
\end{theorem}

\begin{proposition}
\label{lov}
Let $|z|=q^{-\sigma}$ with $1 < \sigma \leq 4/3$ and $i \in \{0,1,\ldots, \ell-1\}$. Further assume that  $|z^{\ell}-q^{-\ell-1}| > \delta$ for some $\delta>0$. For $h \in \mathbb{F}_q[t]$ and $j \geq 1$, we have
\begin{align*}
  \sum_{V \in \mathcal{M}_n} \Big| \psi^{(i)} (hV^j, z ) \Big|^2 & \ll |h|^{\varepsilon} q^{\varepsilon n} \Bigg( q^n + \Big(q^{n(\frac{1}{3}+\frac{3j}{2}-j\sigma)} |h|^{3/2-\sigma} +  q^{n(\frac{2}{3}+\frac{4j}{3}-j\sigma)} |h|^{4/3-\sigma}\Big) \\
  & \times \max \Big\{1, q^{jn(\sigma-7/6)} |h|^{\sigma-7/6} \Big\}\Bigg) .
  \end{align*}
\end{proposition}

 \begin{align}
& \sum_{v \in \mathcal{M}_{n - \deg(\text{rad}(Q))}}  \Bigg| \sum_{\substack{F' \in \mathcal{H} \\  \deg(F') \leq [(\deg(h)+jn)/2]-\deg(ab^2) \nonumber  \\ (F',hQv)=1}} a_{T_1}(F') \overline{\Big( \frac{F'}{v}\Big)_{\ell}} \Bigg|^2\\
 &  \ll |h|^{\varepsilon} q^{\varepsilon n} \Bigg(q^{n-\deg(\text{rad}(Q))}+q^{n(1/3+j/2)} \frac{|h|^{1/2}}{|\text{rad}(Q)|^{1/3} |ab^2|}+q^{n(2/3+j/3)} \frac{|h|^{1/3}}{|\text{rad}(Q)|^{2/3}|ab^2|^{2/3}} \Bigg)   \nonumber \\
 & \times q^{jn(1-\sigma)} |h|^{1-\sigma} |ab^2|^{2\sigma-2} \nonumber \\
 & \ll |h|^{\varepsilon} q^{\varepsilon n} \Bigg(q^{n(1+j-j\sigma)-\deg(\text{rad}(Q))} |h|^{1-\sigma} |ab^2|^{2\sigma-2} + q^{n (1/3+3j/2-j \sigma)} \frac{|h|^{3/2-\sigma}}{|\text{rad}(Q)|^{1/3} |ab^2|^{3-2\sigma}} \nonumber\\
 & +q^{n(2/3+4j/3-j \sigma)} \frac{|h|^{4/3-\sigma}}{|\text{rad}(Q)|^{2/3}|ab^2|^{8/3-2\sigma}}\Bigg). \label{sumv}
\end{align} 

\begin{corollary}
\label{cor_lov}
For $i \in \{0,\ldots,\ell-1\}$, $j \geq 2$ and $z$ as in Proposition \ref{lov}, we have
\begin{align*}
 \sum_{V \in \mathcal{M}_n} \Big|  \psi^{(i)}(hV^j,z)\Big| \ll |h|^{\varepsilon} q^{\varepsilon n} \Big(q^n + \max \{ q^{n(\frac{2}{3} + \frac{3j}{4} - \frac{j \sigma}{2})} |h|^{\frac{1}{2} (\frac{3}{2}-\sigma)}, q^{n (\frac{2}{3}+\frac{j}{6})} |h|^{\frac{1}{6}}\} \Big).
 \end{align*}
\end{corollary}  

 \begin{proposition}
Let $\ell \geq 5$, $i \in \{0,\ldots, \ell-1\}$, and $z$ as in Proposition  \ref{lov}. We have
\begin{align*}
\sum_{V \in \mathcal{M}_n} & |\psi^{(i)}_V(V,b,z)|   \ll |b|^{1/2-\sigma+\varepsilon} q^{\varepsilon n} \Big(q^n +q^{ n \max\{ \frac{\ell(9-6\sigma)-6\sigma-1}{12} , \frac{7-2\sigma}{4} \}} |b|^{(\ell-2)\frac{1}{2} (\frac{3}{2}-\sigma)}\\
& +q^{n \max\{ \frac{\ell}{6}-2\sigma+\frac{5}{3}, \frac{7}{6} \}} |b|^{\frac{\ell-2}{6}} \sum_{n_1+\ldots+\ell n_{\ell}=n}  \prod_{i \geq \max\{2, \frac{ \frac{\ell+1}{6} -\sigma}{\sigma-\frac{2}{3}} \}} q^{n_i (1-\frac{i}{2} ( \frac{\ell}{6} +\frac{15}{6}-3\sigma))}  \Big).
\end{align*}



 \end{proposition}

 \begin{proof}
 Using Theorem \ref{thm:gettingridofa}, we have 
 \begin{align*}
  \psi_V^{(i)}(V, b,z) =&
 G_\ell(r,b) z^{\deg(b)}\sum_{\substack{E\in \mathcal{M}\\ E \mid V_\ell^* }}\mu(E)G_\ell(b^{\ell-2}V_1\cdots V_{\ell-1}^{\ell-1},E)z^{\deg(E)}\\&\times
  \prod_{\pi \mid bEV_1\cdots V_{\ell-1}}
\left(1 - q^{(\ell-1)\deg(\pi)} \legl{-1}{\pi}^{\nu_\pi(aEr_1\cdots V_{\ell-2}V_{\ell-1})+1}z^{\ell \deg(\pi)}\right)^{-1}\nonumber \\
&\times  \sum_{\substack{d_{\ell-2} \mid bEV_{\ell-2}\\ d_1\mid V_1, \dots d_{\ell-3}\mid V_{\ell-3}}} \mu(d_1d_2\cdots d_{\ell-2} ) q^{\deg(d_1d_2^2\cdots d_{\ell-2}^{\ell-2})}  z^{\deg(d_1^2d_2^3\cdots d_{\ell-2}^{\ell-1}  )} \prod_{1\leq k< j \leq \ell-2} \overline{\legl{d_k}{d_j}^{k(j+1)}}\\
& \times \prod_{j=1}^{\ell-2} G_{\ell,j+1} \left( \frac{b^{\ell-2}E^{\ell-2}V_1 V_2^2 \cdots V_{\ell-1}^{\ell-1}}{d_{1}d_2\cdots d_{\ell-2} ^{\ell-2}} , d_{j} \right)\\
&  \times \psi^{(i-\deg(b)-\deg(E)- \deg(d_1^2d_2^3\cdots d_{\ell-2}^{\ell-1} ))} \left( \frac{b^{\ell-2}E^{\ell-2}V_1 V_2^2 \cdots V_{\ell-1}^{\ell-1}}{\prod_{j=1}^{\ell-2}d_j^{2j+2-\ell}}, z \right),
\end{align*}
where we recall that $V=V_1 \ldots V_{\ell-1}^{\ell-1} V_{\ell}^{\ell}$ with $V_1, \ldots, V_{\ell-1}$ square-free and pairwise coprime, and where $V_{\ell}^{*}$ denotes the product of the primes dividing $V_{\ell}$, but not $V_1 \ldots V_{\ell-1}$. It then follows that for any $\varepsilon>0$
\begin{align*}
\sum_{V \in \mathcal{M}_n} & |\psi^{(i)}_V(V,b,z)| \ll |b|^{1/2-\sigma+\varepsilon} \sum_{n_1+2n_2+\cdots+\ell n_\ell=n} \sum_{\substack{i=1 \\ \deg(V_i)=n_i} }^{\ell} \sum_{E|V_{\ell}^*} |E|^{1/2-\sigma+\varepsilon}  \nonumber  \\
& \times \sum_{\substack{d_{\ell-2} \mid bEV_{\ell-2}\\ d_i|V_i, i \leq \ell-3}} |d_i|^{i+1/2-(i+1)\sigma} |d_{\ell-2}|^{\ell-3/2-(\ell-1)\sigma} \sum_{j=0}^{\ell-1} \Big| \psi^{(j)}\left( \frac{b^{\ell-2}E^{\ell-2}V_1 V_2^2 \cdots V_{\ell-1}^{\ell-1}}{\prod_{j=1}^{\ell-2}d_j^{2j+2-\ell}}, z \right) \Big|. %\label{sumv0}
\end{align*}
For simplicity, denote
$$ \alpha = \frac{b^{\ell-2}E^{\ell-2} V_2^2 \cdots V_{\ell-1}^{\ell-1}}{\prod_{j=2}^{\ell-2}d_j^{2j+2-\ell}}.$$
Focusing on the sum over $V_1$ in the expression above and letting $V_1=d_1a$, we have
\begin{align}
\label{sumalov}
\sum_{\deg(V_1)=n_1} \sum_{d_1|V_1}|d_1|^{\frac{3}{2}-2\sigma} |\psi^{(j)}(\alpha V_1 d_1^{\ell-4},z) | = \sum_{\deg(a) \leq n_1} q^{(n_1-\deg(a))(\frac{3}{2}-2\sigma)} \sum_{\deg(d_1)=n_1-\deg(a)} |\psi^{(j)}(\alpha a d_1^{\ell-3} ,z)|.
\end{align}
Now we use Corollary \ref{cor_lov} to get that 
\begin{align*}
\sum_{\deg(d_1)=n_1-\deg(a)}&  |\psi^{(j)}(\alpha a d_1^{\ell-3} ,z)| \ll |\alpha a|^{\varepsilon} q^{\varepsilon n_1} \Big( q^{n_1-\deg(a)}  \\
&+ \max \{ q^{(n_1-\deg(a))(\frac{3\ell}{4} - \frac{19}{12} - \frac{(\ell-3)\sigma}{2})} |\alpha a|^{\frac{1}{2}( \frac{3}{2} -\sigma)} , q^{(n_1-\deg(a))(\frac{\ell}{6}+\frac{1}{6})} |\alpha a|^{\frac{1}{6}}\} \Big) .
\end{align*}
Introducing the sum over $a$, we have that
\begin{align*}
\eqref{sumalov} &\ll q^{n_1(\frac{3}{2}-2\sigma)+\varepsilon n} |\alpha|^{\varepsilon} \sum_{\deg(a) \leq n_1} \frac{1}{|a|^{\frac{3}{2}-2\sigma}} \Big( q^{n_1-\deg(a)}  \\
&+ \max \{ q^{(n_1-\deg(a))(\frac{3\ell}{4} - \frac{19}{12} - \frac{(\ell-3)\sigma}{2})} |\alpha a|^{\frac{1}{2}( \frac{3}{2} -\sigma)} , q^{(n_1-\deg(a))(\frac{\ell}{6}+\frac{1}{6})} |\alpha a|^{\frac{1}{6}}\} \Big) \\
& \ll |\alpha|^{\varepsilon} \Big( q^{n_1} + \max\{ q^{n_1 \max\{ \frac{\ell(9-6\sigma)-6\sigma-1}{12} , \frac{7-2\sigma}{4} \}} |\alpha|^{\frac{1}{2}(\frac{3}{2}-\sigma)}, q^{n_1 \max\{ \frac{\ell}{6}-2\sigma+\frac{5}{3}, \frac{7}{6} \}} |\alpha|^{\frac{1}{6}}\} \Big).
\end{align*}
Now we introduce the sum over $V_i$ for $2 \leq i  \leq \ell-3$. For $\theta \in \{ \frac{1}{2}\left(\frac{3}{2}-\sigma\right), \frac{1}{6}\}$  we have that
$$
\sum_{\deg(V_i) =n_i}   \sum_{d_i|V_i} |d_i|^{i+\frac{1}{2}- (i+1)\sigma} \Big| \frac{V_i^i}{d_i^{2i+2-\ell}} \Big|^{\theta}   \ll
\begin{cases}
q^{n_i(1+i\theta)} & \mbox{ if } \theta(\ell-2-2i)+i+\frac{1}{2} -(i+1)\sigma \leq 0 \\
q^{n_i(\theta(\ell-2-i)+i+\frac{3}{2}-(i+1)\sigma)} & \mbox{ otherwise.}
\end{cases}
$$
We similarly have that
\begin{align*}
\sum_{\deg(V_{\ell-2})=n_{\ell-2}} \sum_{d_{\ell-2}|EbV_{\ell-2}} |d_{\ell-2}|^{\ell-\frac{3}{2}-(\ell-1)\sigma} \Big|\frac{ V_{\ell-2}^{\ell-2}}{d_{\ell-2}^{\ell-2}} \Big|^{\theta} \ll q^{n_{\ell-2}(1+\theta (\ell-2))},
\end{align*}
\begin{align*}
\sum_{\deg(V_{\ell-1})=n_{\ell-1}} |V_{\ell-1}|^{(\ell-1)\theta}\ll q^{n_{\ell-1}(1+\theta (\ell-1))},
\end{align*}
\begin{align*}
\sum_{\deg(V_{\ell}) = n_{\ell}} \sum_{E | V_{\ell}^{*}} |E|^{\frac{1}{2}-\sigma+\varepsilon+(\ell-2)\theta} \ll q^{n_{\ell}(\frac{3}{2}-\sigma+(\ell-2)\theta)}.
\end{align*}
Putting all of these together, we get that
\begin{align}
\eqref{sumv} & \ll |b|^{1/2-\sigma+\varepsilon} \sum_{n_1+2n_2+\cdots+\ell n_{\ell}=n} \Big( q^n +\max\{ q^{n_1 \max\{ \frac{\ell(9-6\sigma)-6\sigma-1}{12} , \frac{7-2\sigma}{4} \}} |b^{\ell-2}|^{\frac{1}{2}\left(\frac{3}{2}-\sigma\right)}  \nonumber \\
& \times  \prod_{\substack{2 \leq i < { \frac{(3-2\sigma)(\ell-2)}{2}+1-2\sigma }}} q^{n_i (\frac{\ell(3-2\sigma)}{4}+\frac{i(1-2\sigma)}{4} )}  \prod_{\substack{i \geq
{ \frac{(3-2\sigma)(\ell-2)}{2}+1-2\sigma }}} q^{n_i (1+\frac{i}{2}\left(\frac{3}{2}-\sigma\right))}  \nonumber \\
&\times q^{n_1 \max\{ \frac{\ell}{6}-2\sigma+\frac{5}{3}, \frac{7}{6} \}} |b^{\ell-2}|^{\frac{1}{6}} \prod_{2 \leq i< \frac{ \frac{\ell+1}{6} -\sigma}{\sigma-\frac{2}{3}}} q^{n_i ( \frac{\ell+7+5i}{6}-(i+1)\sigma)} \prod_{i \geq  \frac{ \frac{\ell+1}{6} -\sigma}{\sigma-\frac{2}{3}}} q^{n_i (1+\frac{i}{6})} \} \Big).\label{mixedbound}
\end{align}
Now note that 
\begin{align*}
 \prod_{\substack{2 \leq i < { \frac{(3-2\sigma)(\ell-2)}{2}+1-2\sigma }}} q^{n_i (\frac{\ell(3-2\sigma)}{4}+\frac{i(1-2\sigma)}{4} )} \leq q^{\left(n-n_1-\sum_{i \geq  \frac{(3-2\sigma)(\ell-2)}{2}+1-2\sigma } in_i\right) \left(  \frac{ \ell(3-2\sigma)}{8} - \frac{2\sigma-1}{4}\right)},
\end{align*}
so \begin{align*}
&  q^{n_1 \max\{ \frac{\ell(9-6\sigma)-6\sigma-1}{12} , \frac{7-2\sigma}{4} \}}   \prod_{\substack{2 \leq i < { \frac{(3-2\sigma)(\ell-2)}{2}+1-2\sigma }}} q^{n_i (\frac{\ell(3-2\sigma)}{4}+\frac{i(1-2\sigma)}{4} )}  \prod_{\substack{i \geq
{ \frac{(3-2\sigma)(\ell-2)}{2}+1-2\sigma }}} q^{n_i (1+\frac{i}{2}\left(\frac{3}{2}-\sigma\right))} \\
& \leq q^{n  (  \frac{ \ell(3-2\sigma)}{8} - \frac{2\sigma-1}{4})} q^{n_1 \Big(  \max\{ \frac{\ell(9-6\sigma)-6\sigma-1}{12} , \frac{7-2\sigma}{4} \} - ( \frac{\ell(3-2\sigma)}{8}-\frac{2\sigma-1}{4})\Big)} \prod_{\substack{i \geq
{ \frac{(3-2\sigma)(\ell-2)}{2}+1-2\sigma }}} q^{n_i \Big(\frac{i\ell(2\sigma-3)}{8}+\frac{i}{2}+1\Big)} \\
& \ll 
q^{ n \max\{ \frac{\ell(9-6\sigma)-6\sigma-1}{12} , \frac{7-2\sigma}{4} \}}.
\end{align*}

We similarly get that
\begin{align*}
& q^{n_1 \max\{ \frac{\ell}{6}-2\sigma+\frac{5}{3}, \frac{7}{6} \}}  \prod_{2 \leq i< \frac{ \frac{\ell+1}{6} -\sigma}{\sigma-\frac{2}{3}}} q^{n_i ( \frac{\ell+7+5i}{6}-(i+1)\sigma)} \prod_{i \geq  \frac{ \frac{\ell+1}{6} -\sigma}{\sigma-\frac{2}{3}}} q^{n_i (1+\frac{i}{6})} \\
& \ll q^{n \max\{ \frac{\ell}{6}-2\sigma+\frac{5}{3}, \frac{7}{6} \}} \prod_{i \geq  \frac{ \frac{\ell+1}{6} -\sigma}{\sigma-\frac{2}{3}}} q^{n_i (1-\frac{i}{2} ( \frac{\ell}{6} +\frac{15}{6}-3\sigma))}.
\end{align*}

Using the two bounds above and \eqref{mixedbound}, we get that
\begin{align*}
\eqref{sumv} & \ll |b|^{1/2-\sigma+\varepsilon} q^{\varepsilon n} \Big(q^n +q^{ n \max\{ \frac{\ell(9-6\sigma)-6\sigma-1}{12} , \frac{7-2\sigma}{4} \}} |b|^{(\ell-2)\frac{1}{2} (\frac{3}{2}-\sigma)}\\
& +q^{n \max\{ \frac{\ell}{6}-2\sigma+\frac{5}{3}, \frac{7}{6} \}} |b|^{\frac{\ell-2}{6}} \sum_{n_1+2n_2+\cdots+\ell n_{\ell}=n}  \prod_{i \geq  \frac{ \frac{\ell+1}{6} -\sigma}{\sigma-\frac{2}{3}}} q^{n_i (1-\frac{i}{2} ( \frac{\ell}{6} +\frac{15}{6}-3\sigma))}  \Big) \Big),
\end{align*}
which finishes the proof.
 \end{proof}

\end{document}
